\cotefile{156-7}
\scanpage{94}
\mspage{93}

J'aimerais prouver le

\uline{\unsure{Théorème}} \textcircled{a} L'application \unsure{précédente} est \unsure{injective}, plus précisément elle induit un \uline{iso. d'ens ordonnés} de $P(\widetilde F)$ sur \deleted{\unsure{le}} un sous-ensemble ordonné de $\Sigma_M$. \textcircled{b} Cette application \deleted{\unsure{est}} commute aux $\operatorname{Sup}$ \margin{facile. $\checkmark$} (tautologiquement) et aux $\operatorname{Inf}$ quelconques. \textcircled{c} On a, pour deux $A,\deleted{\unsure{A}}B\in P(\widetilde F)$, avec $B\subset A$
\[
\deleted{\unsure{S_{B\setminus A}|S_{B\cap A}}=\unsure{S_B}}
\qquad S_{A\setminus B}=S_A\cap C_M(S_B)
\]
\[
\deleted{\unsure{\operatorname{Supp}F}=\unsure{\operatorname{Supp}|S_{\widetilde F}}}
\]
où $C_M$ est l'anti-involution canonique de l'ens. des supports. (\unsure{application cosupport}).

\textcircled{d} \deleted{On a} Si $A$ est une partie \uuline{fermée} de $\widetilde F$ (donc de la forme $\widetilde{\unsure{G|C}}$, où $G\subseteq F$ est une sous-figure de $\unsure{A|\Lambda}$) on a
\[
S_A=\operatorname{Supp}G\overset{\mathrm{def}}=\operatorname{Supp}_{\widehat M}(A)
\]

\uline{NB} En fait, $S_A$ est défini \unsure{sans référence} à la figure ambiante \deleted{\unsure{\(\Phi\)}} $F$, où $\widetilde F=\Phi\subset M$, en termes de l'ens $A\subset M$ formé des multistrates mutuellement compatibles. Donc \unsure{tout} ce formulaire précédent peut se
